\Needspace{0.9\textheight}
\section{Kiến trúc tổng thể và tổ chức tri thức}
\label{sec:kien-truc}
\subsection{Lựa chọn hướng giải quyết}
Các yêu cầu ở Mục~\ref{sec:yeu-cau} đòi hỏi hệ thống vừa hiểu cách diễn đạt
của người hỏi, vừa sử dụng được căn cứ và duy trì thông tin qua nhiều lượt.
Bảng~\ref{tab:so-sanh} đối chiếu những hướng giải quyết theo vai trò của chúng
đối với bài toán; các hướng có thể được kết hợp trong cùng một hệ thống.

\begin{table}[H]
\centering\small
\begin{tabularx}{\textwidth}{|P{27mm}|Y|Y|}
\hline
\textbf{Hướng} & \textbf{Vai trò phù hợp} & \textbf{Phần cần bổ sung cho bài toán} \\\hline
Tra cứu, hỏi đáp thường gặp & Cung cấp nguồn hoặc đáp án đã biên soạn theo chủ đề. & Chọn hướng dẫn theo dữ kiện và theo dõi thay đổi của tình huống. \\ \hline
Dựa trên luật & Đối chiếu điều kiện rõ ràng và truy vết căn cứ của kết luận. & Tiếp nhận cách diễn đạt tự do; duy trì quy tắc khi nguồn thay đổi. \\ \hline
Mô hình ngôn ngữ & Hiểu câu hỏi và diễn đạt phản hồi bằng ngôn ngữ tự nhiên. & Ràng buộc câu trả lời bằng nguồn và điều kiện áp dụng. \\ \hline
Hỏi đáp có truy xuất & Đưa đoạn tài liệu liên quan vào quá trình tạo câu trả lời. & Kiểm tra đối tượng, thời điểm và cập nhật kết luận khi dữ kiện đổi. \\
\hline
\end{tabularx}
\caption{Đối chiếu các hướng giải quyết theo yêu cầu tư vấn}
\label{tab:so-sanh}
\end{table}

Phương án đề xuất phân chia trách nhiệm giữa mô hình ngôn ngữ, truy xuất,
kiểm tra điều kiện và quản lý ngữ cảnh. Mô hình xử lý ngôn ngữ; truy xuất tìm
căn cứ; điều kiện đã xác nhận ràng buộc kết luận; ngữ cảnh duy trì tình huống
đang xét. Sự phối hợp này là cơ sở thiết kế để đáp ứng F01--F05; mức đáp ứng
thực tế được kiểm tra theo định hướng tại Mục~\ref{sec:kiem-chung}.

\Needspace{18\baselineskip}
\subsection{Kiến trúc tổng thể}
Hình~\ref{fig:kien-truc} mô tả các thành phần và quan hệ trao đổi thông tin.
Giao diện tiếp nhận câu hỏi; trạng thái phiên giữ thông tin hiện tại; kho
văn bản và thẻ tư vấn cung cấp căn cứ cho bước đối chiếu. Bộ tạo phản hồi
nhận các kết luận cùng nguồn hỗ trợ để diễn đạt thành hướng dẫn.

\begin{figure}[H]
\centering\begin{tikzpicture}[node distance=9mm and 12mm]
\node[block,text width=35mm] (ui) {Giao diện hội thoại};
\node[block,text width=43mm,right=of ui] (state) {Trạng thái phiên\\Dữ kiện và thay đổi};
\node[block,text width=35mm,right=of state] (nlp) {Mô hình ngôn ngữ\\Hiểu và diễn đạt};
\node[block,text width=35mm,below=of ui] (retrieve) {Truy xuất nguồn\\Lọc quyền và phạm vi};
\node[block,text width=43mm,right=of retrieve] (check) {Kiểm tra điều kiện\\và phụ thuộc kết luận};
\node[block,text width=35mm,right=of check] (compose) {Tạo câu trả lời\\Có bằng chứng};
\node[block,text width=35mm,below=of retrieve] (store) {Kho văn bản\\Phiên bản và đoạn nguồn};
\node[block,text width=43mm,right=of store] (cards) {Thẻ tư vấn\\Điều kiện đã xác nhận};
\draw[arrow] (ui)--(state);
\draw[{Stealth[length=2mm]}-{Stealth[length=2mm]},line width=0.55pt] (state)--(nlp);
\draw[arrow] (state)--(check);
\draw[{Stealth[length=2mm]}-{Stealth[length=2mm]},line width=0.55pt] (nlp)--(compose);
\draw[arrow] (store)--(retrieve);
\draw[arrow] (retrieve)--(check);
\draw[arrow] (cards)--(check);
\draw[arrow] (check)--(compose);
\draw[arrow] (state.south west)--(retrieve.north east);
\end{tikzpicture}

\caption{Kiến trúc tổng thể của hệ thống đề xuất}
\label{fig:kien-truc}
\end{figure}

Các mô hình ngôn ngữ dựa trên kiến trúc Transformer~\cite{vaswani2017attention}
cung cấp khả năng xử lý ngữ cảnh và tạo văn bản. Trong hệ thống, mô hình
ngôn ngữ lớn (LLM) được dùng để nhận diện mục đích hỏi, trích xuất dữ kiện
ứng viên và diễn đạt phản hồi. Thông tin còn mơ hồ cần được xác nhận lại;
mô hình không tự suy ra hồ sơ cá nhân hoặc đặt điều kiện nghiệp vụ.

RAG kết hợp truy xuất tài liệu với mô hình sinh~\cite{lewis2020rag}. Thành
phần truy xuất đưa các đoạn căn cứ liên quan vào ngữ cảnh tạo câu trả lời,
đồng thời lọc theo quyền và phạm vi sử dụng. Các phương pháp mã hóa câu hỏi
và đoạn văn như Dense Passage Retrieval (DPR)~\cite{karpukhin2020dpr} cung cấp cơ sở cho tìm kiếm theo
ngữ nghĩa. Đối với đề tài, đoạn có nội dung gần câu hỏi vẫn phải được kiểm
tra điều kiện áp dụng trước khi dùng để đưa ra kết luận.

Bộ kiểm tra điều kiện sử dụng dữ kiện hiện tại và những quy tắc đã được
xác nhận từ nguồn. Trạng thái phiên liên kết dữ kiện với căn cứ và hướng
dẫn đang dùng, phục vụ việc xét lại khi người hỏi thay đổi thông tin. Cách
phối hợp các thành phần trong một lượt hỏi đáp được trình bày ở
Mục~\ref{sec:co-che}.

\subsection{Tổ chức tri thức tư vấn}
\label{subsec:tri-thuc}
Tri thức được tổ chức thành văn bản gốc, đoạn căn cứ, điều kiện và thẻ tư
vấn. Đoạn căn cứ giữ vị trí điều, khoản, trang hoặc bảng để kiểm tra nội
dung. Điều kiện biểu diễn yêu cầu cần đối chiếu với dữ kiện. Thẻ tư vấn
tập hợp mục đích hỏi, phạm vi, các dữ kiện quyết định, điều kiện liên quan
và đầu mối cần trao đổi; đây là đơn vị tổ chức nội dung, không phải đáp án
cố định hay biểu mẫu bắt buộc người dùng điền.

\begin{table}[H]
\centering\small
\begin{tabularx}{\textwidth}{|P{30mm}|Y|}
\hline
\textbf{Thành phần} & \textbf{Thông tin cần lưu} \\\hline
Nguồn và đoạn căn cứ & Xuất xứ, vị trí trích dẫn, phiên bản, hiệu lực, đối tượng, đơn vị, sửa đổi, chuyển tiếp và quyền sử dụng. \\ \hline
Điều kiện & Nội dung cần kiểm tra, dữ kiện có ảnh hưởng, phạm vi áp dụng và đoạn nguồn xác lập điều kiện. \\ \hline
Thẻ tư vấn & Mục đích, phạm vi hỗ trợ, dữ kiện cần làm rõ, điều kiện liên quan, điểm dừng và chủ nội dung. \\
\hline
\end{tabularx}
\caption{Các thành phần thông tin của kho tri thức tư vấn}
\label{tab:du-lieu}
\end{table}

Quá trình chuẩn bị nguồn gồm thu thập bản gốc, kiểm tra độ đầy đủ, xác minh
phiên bản và phạm vi, trích đoạn, xây điều kiện và xác nhận nội dung. Khi
chia đoạn, cần giữ liên kết tới định nghĩa, ngoại lệ, phụ lục và văn bản dẫn
chiếu có ảnh hưởng. Một điều kiện có thể được nhiều thẻ tham chiếu để hạn
chế sao chép cùng quy tắc thành các bản khác nhau.

Khi nguồn thay đổi, những thẻ và điều kiện liên quan được xác định để rà
soát. Phiên bản cũ được giữ khi còn cần xét chuyển tiếp; tài liệu chưa xác
minh được quản lý riêng. Kho nguồn và trạng thái xác nhận nội dung là điều
kiện để lựa chọn căn cứ phù hợp cho mỗi tình huống.
